\corrigesChapitre{17}{Puissance et énergie électrique}
\par\noindent\begin{minipage}{\linewidth}
\begin{corrige}{1}{Deux grandeurs}
\begin{questions}
\item Une puissance.
\item J, Wh ou kWh, selon le contexte.
\end{questions}
\end{corrige}
\end{minipage}
\par\noindent\begin{minipage}{\linewidth}
\begin{corrige}{2}{Une heure}
\begin{questions}
\item $5\unite{Wh}$.
\item $5\times3\,600=18\,000\unite J$.
\end{questions}
\par\smallskip\noindent\textbf{Pour comprendre.}\enspace À puissance constante, $E=P\Delta t=5\times1=5\unite{Wh}$. Puis $1\unite{Wh}=3\,600\unite J$ : $E=5\times3\,600=18\,000\unite J$.\par
\end{corrige}
\end{minipage}
\par\noindent\begin{minipage}{\linewidth}
\begin{corrige}{3}{Convertir}
\begin{questions}
\item $1\,500\unite{Wh}$.
\item $0{,}5\unite h$.
\end{questions}
\end{corrige}
\end{minipage}
\par\noindent\begin{minipage}{\linewidth}
\begin{corrige}{4}{Un produit électrique}
\begin{questions}
\item $P=UI=6\times0{,}20=1{,}2\unite W$.
\item Le watt.
\end{questions}
\end{corrige}
\end{minipage}
\par\noindent\begin{minipage}{\linewidth}
\begin{corrige}{5}{La durée compte}
\begin{questions}
\item A : 20 Wh contre 10 Wh.
\item Oui : à durée égale, l’énergie est proportionnelle à la puissance.
\end{questions}
\end{corrige}
\end{minipage}
\par\noindent\begin{minipage}{\linewidth}
\begin{corrige}{6}{En secondes}
\begin{questions}
\item $E=3\times40=120\unite J$.
\item La durée est en secondes ; W multiplié par s donne J.
\end{questions}
\end{corrige}
\end{minipage}
\par\noindent\begin{minipage}{\linewidth}
\begin{corrige}{7}{Une demi-heure}
\begin{questions}
\item $E=60\times0{,}5=30\unite{Wh}=0{,}030\unite{kWh}$.
\item $30\times3\,600=108\,000\unite J$.
\end{questions}
\par\smallskip\noindent\textbf{Pour comprendre.}\enspace Convertir d’abord 30 min en 0,50 h : $E=60\times0{,}50=30\unite{Wh}=0{,}030\unite{kWh}$. En SI, 30 min = 1 800 s et $E=60\times1\,800=108\,000\unite J$.\par
\end{corrige}
\end{minipage}
\par\noindent\begin{minipage}{\linewidth}
\begin{corrige}{8}{Retrouver une durée}
\begin{questions}
\item $\Delta t=E/P=180/6=30\unite s$.
\item $6\times30=180\unite J$.
\end{questions}
\par\smallskip\noindent\textbf{Pour comprendre.}\enspace Isoler la durée : $\Delta t=E/P=180/6=30\unite s$. L’unité J/W donne des secondes ; le contrôle $6\times30=180\unite J$ retrouve l’énergie reçue.\par
\end{corrige}
\end{minipage}
\par\noindent\begin{minipage}{\linewidth}
\begin{corrige}{9}{Une intensité nominale}
\begin{questions}
\item $I=P/U=2{,}4/6=0{,}40\unite A$.
\item $400\unite{mA}$.
\end{questions}
\end{corrige}
\end{minipage}
\par\noindent\begin{minipage}{\linewidth}
\begin{corrige}{10}{Un coût fictif}
\begin{questions}
\item $2{,}5\times0{,}20=0{,}50$ euro.
\item Non : le calcul ne porte que sur le poste défini dans l’énoncé.
\end{questions}
\end{corrige}
\end{minipage}
\par\noindent\begin{minipage}{\linewidth}
\begin{corrige}{11}{Puissante mais brève}
\begin{questions}
\item A : $100\times60=6\,000\unite J$ ; B : $10\times1\,200=12\,000\unite J$.
\item B reçoit deux fois plus d’énergie malgré sa puissance plus faible, car sa durée est beaucoup plus longue.
\end{questions}
\par\smallskip\noindent\textbf{Pour comprendre.}\enspace Convertir les durées avant de comparer : 1 min = 60 s et 20 min = 1 200 s. A reçoit 6 000 J, B 12 000 J. La plus grande puissance ne suffit donc pas à désigner la plus grande énergie.\par
\pcFigureEx{figures/cor-puissance-duree.pdf}{Comparer des énergies exige de tenir compte à la fois de la puissance et de la durée.}
\end{corrige}
\end{minipage}
\par\noindent\begin{minipage}{\linewidth}
\begin{corrige}{12}{Une veille}
\begin{questions}
\item $40\unite{Wh}$ par jour ; $1\,200\unite{Wh}=1{,}2\unite{kWh}$ pour 30 jours.
\item Les autres périodes de fonctionnement et éventuellement les postes inclus dans un périmètre plus large.
\end{questions}
\end{corrige}
\end{minipage}
\par\noindent\begin{minipage}{\linewidth}
\begin{corrige}{13}{Une résistance qui chauffe}
\begin{questions}
\item $I=5/100=0{,}050\unite A$ ; $P=5\times0{,}050=0{,}25\unite W$.
\item Non : $0{,}25>0{,}125\unite W$. Le composant serait sollicité au-delà de sa valeur admissible.
\end{questions}
\par\smallskip\noindent\textbf{Pour comprendre.}\enspace Calculer successivement $I=U/R=5/100=0{,}050\unite A$, puis $P=UI=5\times0{,}050=0{,}25\unite W$. La limite 0,125 W est dépassée d’un facteur deux : ne pas réaliser ce fonctionnement.\par
\end{corrige}
\end{minipage}
\par\noindent\begin{minipage}{\linewidth}
\begin{corrige}{14}{Même service}
\begin{questions}
\item A : $4\unite{kWh}$ ; B : $0{,}8\unite{kWh}$ ; économie : $3{,}2\unite{kWh}$.
\item Vérifier que qualité et quantité de lumière sont comparables ; ce calcul ne comprend pas tout le cycle de vie.
\end{questions}
\end{corrige}
\end{minipage}
\par\noindent\begin{minipage}{\linewidth}
\begin{corrige}{15}{Une unité fausse}
\begin{questions}
\item kW/h n’est pas une unité d’énergie ; il faut multiplier puissance et durée.
\item $E=2\times2=4\unite{kWh}$.
\end{questions}
\end{corrige}
\end{minipage}
\par\noindent\begin{minipage}{\linewidth}
\phantomsection\label{cor-defi-17}\pcSolutionDefi{17}{Le projet d’éclairage}\begin{questions}
\item $E=10\times12\times4\times20=9\,600\unite{Wh}=9{,}6\unite{kWh}$.
\item Nouvelle énergie : $6{,}4\unite{kWh}$ ; économie : $3{,}2\unite{kWh}$.
\item $3{,}2\times0{,}25=0{,}80$ euro. Ce gain ne décide pas seul d’un achat : prix, durée de vie et impacts doivent être considérés.
\end{questions}
\end{minipage}
\par\noindent\begin{minipage}{\linewidth}
\phantomsection\label{cor-activite-17}\pcSousTitre{Pistes pour l’investigation}
À même durée, multiplier chaque puissance par cette durée ; convertir Wh en kWh en divisant par 1 000. Si les durées diffèrent, les produits doivent être recalculés avant de comparer. Une puissance plus faible ne garantit pas une énergie totale plus faible. Le service rendu et le périmètre de l’étude conditionnent la conclusion.\par\end{minipage}
